\ifdefined\gkver\else\def\gkver{student}\fi
\documentclass[\gkver]{gaokaozhenti}
\begin{document}

\chapter{2010年上海理综}
%% paperType: 理综物理部分

\begin{enumerate}
\item
%% number: 4
%% paperName: 2010年上海理综第 4 题 3 分
%% typeId: 1
%% type: 单选题
%% chapter: 宇宙恒星演化
%% point: 宇宙的基本结构
%% method: 无
%% score: 3
%% degree: 850
%% duplicateId: 0
%% body:
右图是一张天文爱好者经过长时间拍摄的“星星的轨道”照片。这些规律的弧线的形成，说明了\xzanswer{D}

\onepicture[w=4.2cm]{figs/04.png}

\fourchoices[answer=D]
{太阳在运动}
{月球在运动}
{地球在公转}
{地球在自转}

%% answer: D
%% memo:
\memoanswer{
本题原卷及材料中未提供详解，待补充。
}

\item
%% number: 5
%% paperName: 2010年上海理综第 5 题 3 分
%% typeId: 1
%% type: 单选题
%% chapter: 功和能
%% point: 机械能守恒定律
%% method: 无
%% score: 3
%% degree: 950
%% duplicateId: 0
%% body:
高台滑雪运动员腾空跃下，如果不考虑空气阻力，则下落过程中该运动员机械能的转换关系是\xzanswer{C}

\fourchoices[answer=C]
{动能减少，重力势能减少}
{动能减少，重力势能增加}
{动能增加，重力势能减少}
{动能增加，重力势能增加}

%% answer: C
%% memo:
\memoanswer{
本题原卷及材料中未提供详解，待补充。
}

\item
%% number: 6
%% paperName: 2010年上海理综第 6 题 3 分
%% typeId: 1
%% type: 单选题
%% chapter: 气体性质
%% point: 查理定律
%% method: 无
%% score: 3
%% degree: 750
%% duplicateId: 0
%% body:
民间常用“拔火罐”来治疗某些疾病，方法是将点燃的纸片放入一个小罐内，当纸片燃烧完时，迅速将火罐开口端紧压在皮肤上，火罐就会紧紧地被“吸”在皮肤上。其原因是，当火罐内的气体\xzanswer{B}

\fourchoices[answer=B]
{温度不变时，体积减小，压强增大}
{体积不变时，温度降低，压强减小}
{压强不变时，温度降低，体积减小}
{质量不变时，压强增大，体积减小}

%% answer: B
%% memo:
\memoanswer{
本题原卷及材料中未提供详解，待补充。
}

\item
%% number: 7
%% paperName: 2010年上海理综第 7 题 3 分
%% typeId: 1
%% type: 单选题
%% chapter: 磁场
%% point: 左手定则
%% method: 无
%% score: 3
%% degree: 750
%% duplicateId: 0
%% body:
如右图，水平桌面上放置一根条形磁铁，磁铁中央正上方用绝缘弹簧悬挂一水平直导线，并与磁铁垂直。当直导线中通入图中所示方向的电流时，可以判断出\xzanswer{A}

\onepicture[w=5cm]{figs/07.png}

\fourchoices[answer=A]
{弹簧的拉力增大，条形磁铁对桌面的压力减小}
{弹簧的拉力减小，条形磁铁对桌面的压力减小}
{弹簧的拉力增大，条形磁铁对桌面的压力增大}
{弹簧的拉力减小，条形磁铁对桌面的压力增大}

%% answer: A
%% memo:
\memoanswer{
本题原卷及材料中未提供详解，待补充。
}

\item
%% number: 8
%% paperName: 2010年上海理综第 8 题 3 分
%% typeId: 1
%% type: 单选题
%% chapter: 曲线运动
%% point: 线速度角速度
%% method: 无
%% score: 3
%% degree: 750
%% duplicateId: 0
%% body:
右图是位于锦江乐园的摩天轮，高度为 $108\Um$，直径是 $98\Um$。一质量为 $50\Ukg$ 的游客乘坐该摩天轮做匀速圆周运动旋转一圈需 $25\Umin$。如果以地面为零势能面，则他到达最高处时的（取 $g=10\Umsq$）\xzanswer{C}

\onepicture[w=4.2cm]{figs/08.jpg}

\fourchoices[answer=C]
{重力势能为 $5.4\times10^{4}\UJ$，角速度为 $0.2\Urads$}
{重力势能为 $4.9\times10^{4}\UJ$，角速度为 $0.2\Urads$}
{重力势能为 $5.4\times10^{4}\UJ$，角速度为 $4.2\times10^{-3}\Urads$}
{重力势能为 $4.9\times10^{4}\UJ$，角速度为 $4.2\times10^{-3}\Urads$}

%% answer: C
%% memo:
\memoanswer{
本题原卷及材料中未提供详解，待补充。
}

\item
%% number: 9
%% paperName: 2010年上海理综第 9 题 3 分
%% typeId: 1
%% type: 单选题
%% chapter: 光学
%% point: 光的电磁说
%% method: 无
%% score: 3
%% degree: 550
%% duplicateId: 0
%% body:
使用照相机拍摄清晰满意的照片，必须选择合适的曝光量。曝光量 $P$ 可以表示为：$P=k\left(\frac{d}{f}\right)^{2}t$，式中 $k$ 为常数，$d$ 为照相机镜头“通光孔径”的直径，$f$ 为照相机镜头的焦距，$t$ 为曝光时间，将 $\frac{d}{f}$ 的倒数 $\frac{f}{d}$ 称为照相机的光圈数。一摄影爱好者在某次拍摄时，选择的光圈数是 8，曝光时间是 $\frac{1}{30}\Us$。若他想把曝光时间减少一半，但不改变曝光量，那么光圈数应选择\xzanswer{B}

\fourchoices[answer=B]
{4}
{5.6}
{8}
{11}

%% answer: B
%% memo:
\memoanswer{
本题原卷及材料中未提供详解，待补充。
}

\item
%% number: 40
%% paperName: 2010年上海理综第 40 题 2 分
%% typeId: 3
%% type: 填空题
%% chapter: 气体性质
%% point: 能源的开发与利用
%% method: 无
%% score: 2
%% degree: 850
%% duplicateId: 0
%% body:
太阳能是清洁的新能源，为了环保，我们要减少使用像煤炭这样的常规能源而大力开发新能源。划分下列能源的类别，把编号填入相应的表格。

①石油　②地热能　③核能　④天然气　⑤风能　⑥潮汐能

\begin{table}[htbp]
\centering
\begin{tblr}{|c|c|}
\hline
类别 & 编号 \\
\hline
新能源 & \tkanswer[2cm]{②③⑤⑥} \\
\hline
常规能源 & \tkanswer[2cm]{①④} \\
\hline
\end{tblr}
\end{table}

%% answer: 
\jdanswer{
②③⑤⑥；①④
}
%% memo:
\memoanswer{
本题原卷及材料中未提供详解，待补充。
}

\item
%% number: 41
%% paperName: 2010年上海理综第 41 题 2 分
%% typeId: 3
%% type: 填空题
%% chapter: 电路
%% point: 电功电功率
%% method: 无
%% score: 2
%% degree: 850
%% duplicateId: 0
%% body:
中国馆、世博中心和主题馆等主要场馆，太阳能的利用规模达到了历届世博会之最，总发电装机容量达到 $4.6\times10^{3}\UkW$。设太阳能电池板的发电效率为 $18\%$，已知地球表面每平方米接收太阳能的平均辐射功率为 $1.353\UkW$，那么所使用的太阳能电池板的总面积为\tkanswer[2cm]{$1.9\times10^{4}$} $\Umq$。

\onepicture[align=right, w=3.2cm]{figs/41.png}

%% answer: 
\jdanswer{
$1.9\times10^{4}$
}
%% memo:
\memoanswer{
本题原卷及材料中未提供详解，待补充。
}

\item
%% number: 42
%% paperName: 2010年上海理综第 42 题 3 分
%% typeId: 3
%% type: 填空题
%% chapter: 电路
%% point: 模块电路
%% method: 无
%% score: 3
%% degree: 850
%% duplicateId: 0
%% body:
各场馆的机器人非常引人注目。在下图设计的机器人模块中，分别填入传感器和逻辑门的名称，使该机器人能够在明亮的条件下，听到呼唤声就来为你服务。

\onepicture[align=right, w=9cm]{figs/42.png}

%% answer: 
\jdanswer{
光；声；与（\&）
}
%% memo:
\memoanswer{
本题原卷及材料中未提供详解，待补充。
}

\item
%% number: 43
%% paperName: 2010年上海理综第 43 题 4 分
%% typeId: 3
%% type: 填空题
%% chapter: 运动和力
%% point: 牛顿第二定律
%% method: 无
%% score: 4
%% degree: 750
%% duplicateId: 0
%% body:
纯电动概念车 E1 是中国馆的镇馆之宝之一。若 E1 概念车的总质量为 $920\Ukg$，在 $16\Us$ 内从静止加速到 $100\Ukmh$（即 $27.8\Ums$），受到恒定的阻力为 $1500\UN$，假设它做匀加速直线运动，其动力系统提供的牵引力为\tkanswer[1.5cm]{$3.1\times10^{3}$} $\UN$。当 E1 概念车以最高时速 $120\Ukmh$（即 $33.3\Ums$）做匀速直线运动时，其动力系统输出的功率为\tkanswer[1.5cm]{50} $\UkW$。

\onepicture[align=right, w=4.5cm]{figs/43.png}

%% answer: 
\jdanswer{
$3.1\times10^{3}$；50
}
%% memo:
\memoanswer{
本题原卷及材料中未提供详解，待补充。
}

\item
%% number: 44
%% paperName: 2010年上海理综第 44 题 5 分
%% typeId: 3
%% type: 填空题
%% chapter: 电路
%% point: 多用表的使用
%% method: 无
%% score: 5
%% degree: 650
%% duplicateId: 0
%% body:
在世博园区，运行着许多氢燃料汽车，其动力来源是氢燃料电池（结构如图）。

\begin{enumerate}
	\item
	以下是估测氢燃料电池输出功率的实验步骤：
	\begin{steps}
		\item
		把多用表的选择开关调至电流档，并选择恰当量程，串联在电路中。读出电流 $I$；
		\item
		把多用表的选择开关调至电压档，把红、黑表笔并联在电动机两端，其中红表笔应该接在图中\tkanswer[1cm]{A}（填“A”或“B”）端。读出电压 $U$；
		\item
		重复步骤①和②，多次测量，取平均值；
		\item
		根据公式 $P=$\tkanswer[1.5cm]{$UI$} 计算氢燃料电池输出功率。
	\end{steps}
	\item
	在上述第②步中遗漏的操作是\tkanswer[2.5cm]{选择恰当量程}；
	\item
	如果该电动机的效率为 $\eta$，汽车运动的速度为 $v$，则汽车的牵引力为\tkanswer[2cm]{$\eta UI/v$}。
\end{enumerate}

\onepicture[align=right, w=5.5cm]{figs/44a.png}

%% answer: 
\jdanswer{
\begin{enumerate}
	\item
	A；$UI$

	\item
	选择恰当量程

	\item
	$\eta UI/v$

\end{enumerate}
}
%% memo:
\memoanswer{
本题原卷及材料中未提供详解，待补充。
}

\end{enumerate}

\end{document}
